\documentclass[12pt, a4paper]{ctexart}
\usepackage{fontspec}
\usepackage{xcolor}
\usepackage{hyperref}
\usepackage{geometry}
\usepackage{enumitem}
\usepackage{parskip}
\usepackage{titlesec}
\usepackage{tcolorbox}
\tcbuselibrary{breakable}
\tcbset{autoparskip}

\usepackage{setspace}
\usepackage{listings}
\usepackage{amssymb}

\geometry{left=2.5cm, right=2.5cm, top=2.5cm, bottom=2.5cm}

\setmainfont{Times New Roman}
\setCJKmainfont{SimSun}

\hypersetup{
    colorlinks=true,
    linkcolor=blue!70!black,
    urlcolor=cyan,
    pdftitle={面向对象（下）-逐题精讲解义},
    pdfauthor={专升本-fifine老师}
}

\definecolor{myblue}{RGB}{30,100,200}
\definecolor{lightblue}{RGB}{230,240,255}
\definecolor{darkblue}{RGB}{0,50,120}

\newtcolorbox{bluebox}[1][]{
    breakable,
    colback=lightblue,
    colframe=myblue,
    coltitle=darkblue,
    fonttitle=\bfseries,
    title={#1},
    boxrule=0.8pt,
    arc=4pt,
    left=6pt,
    right=6pt,
    top=6pt,
    bottom=6pt,
    before skip=8pt,
    after skip=8pt
}

\usepackage{etoolbox}
\BeforeBeginEnvironment{bluebox}{\par\vfill}%
\AfterEndEnvironment{bluebox}{\par\vfill}%

\titleformat{\section}
  {\normalfont\Large\bfseries\color{myblue}}{\thesection}{1em}{}
\titlespacing*{\section}{0pt}{3.5ex plus 1ex minus .2ex}{1.5ex plus .2ex}

\title{\textcolor{myblue}{\textbf{面向对象（下） - 逐题精讲解义}}}
\author{专升本-fifine老师 教学组}
\date{2026年03月05日}

\lstset{
    language=Java,
    basicstyle=\ttfamily\small,
    keywordstyle=\color{blue},
    commentstyle=\color{green!60!black},
    stringstyle=\color{orange},
    numbers=left,
    numberstyle=\tiny\color{gray},
    frame=single,
    breaklines=true,
    columns=fullflexible,
    showstringspaces=false
}

\newcommand{\code}[1]{\texttt{#1}}

\begin{document}

\maketitle
\thispagestyle{empty}
\newpage

\section{面向对象（下）}

\begin{enumerate}[label=\textbf{\arabic*.}]

	\item \textbf{在Java中，以下哪个关键字用于实现继承？}
	      \begin{enumerate}[label=\Alph*.]
		      \item extends
		      \item implements
		      \item interface
		      \item abstract
	      \end{enumerate}

	      \begin{bluebox}{答案与深度解析}
		      \textbf{答案：A. extends}

		      \textbf{【核心认知框架】}
		      \begin{itemize}[label={}, leftmargin=*, nosep]
			      \item \textbf{题目本质}：考查类继承与接口实现的关键字区分
			      \item \textbf{关键区分点}：extends继承类，implements实现接口
			      \item \textbf{命题人意图}：测试面向对象核心语法——继承的实现方式
		      \end{itemize}

		      \textbf{【选项原理深度拆解】}

		      \begin{itemize}[leftmargin=*, nosep, label={}]
			      \item[A] \textbf{extends — 正确（类继承关键字）}
			            \begin{itemize}[label={}, leftmargin=*, nosep]
				            \item \textbf{语法规则}：class 子类 extends 父类（JLS §8.1.4）
				            \item \textbf{继承特性}：单继承（一个类只能继承一个父类）
				            \item \textbf{字节码证据}：
				                  \begin{lstlisting}[language=Java]
// 源码
class Animal {}
class Dog extends Animal {}

// 字节码（javap -v）
// class Dog extends Animal
// Super class: #2; //class Animal
				      \end{lstlisting}
				            \item \textbf{内存模型}：
				                  \begin{itemize}[label={}, nosep]
					                  \item Class对象：Dog继承Animal的类信息
					                  \item 方法表：包含继承的方法
					                  \item 实例布局：包含父类+子类字段
				                  \end{itemize}
				            \item \textbf{完整代码示例}：
				                  \begin{lstlisting}[language=Java]
// 父类
class Animal {
    String name;

    public void eat() {
        System.out.println("动物吃东西");
    }
}

// 使用extends继承父类
class Dog extends Animal {
    String breed;  // 子类特有字段

    @Override
    public void eat() {
        System.out.println("狗吃骨头");  // 重写父类方法
    }

    public void bark() {
        System.out.println("汪汪叫");  // 子类特有方法
    }
}

public class ExtendsDemo {
    public static void main(String[] args) {
        Dog dog = new Dog();
        dog.name = "旺财";         // 继承的属性
        dog.breed = "金毛";        // 子类属性
        dog.eat();                 // 继承的方法（被重写）
        dog.bark();                // 子类方法
    }
}
				            \end{lstlisting}
				            \item \textbf{继承层次}：
				                  \begin{lstlisting}[language=Java]
// 继承链：Object -> Animal -> Dog
class Dog extends Animal {}       // 直接继承
class Puppy extends Dog {}        // 间接继承Animal
// 多层继承支持
				            \end{lstlisting}
			            \end{itemize}

			      \item[B] \textbf{implements — 错误（接口实现关键字）}
			            \begin{itemize}[label={}, leftmargin=*, nosep]
				            \item \textbf{语法规则}：class 类名 implements 接口1, 接口2（JLS §8.1.5）
				            \item \textbf{特性}：多实现（一个类可实现多个接口）
				            \item \textbf{字节码证据}：
				                  \begin{lstlisting}[language=Java]
// 源码
interface Flyable { void fly(); }
class Bird implements Flyable {
    public void fly() { System.out.println("飞翔"); }
}

// 字节码
// interfaces: 2, fields: 0, methods: 1
// 接口表包含Flyable
				      \end{lstlisting}
				            \item \textbf{代码示例}：
				                  \begin{lstlisting}[language=Java]
interface Swimable { void swim(); }
interface Flyable { void fly(); }

class Duck implements Swimable, Flyable {
    @Override
    public void swim() {
        System.out.println("鸭子游泳");
    }

    @Override
    public void fly() {
        System.out.println("鸭子飞翔");
    }
}
// 一个类可实现多个接口
				      \end{lstlisting}
				            \item \textbf{命题陷阱}：考生容易混淆类继承和接口实现
			            \end{itemize}

			      \item[C] \textbf{interface — 错误（接口定义关键字）}
			            \begin{itemize}[label={}, leftmargin=*, nosep]
				            \item \textbf{语法规则}：interface 接口名（JLS §9.1）
				            \item \textbf{作用}：定义接口规范，不是实现继承
				            \item \textbf{字节码特征}：
				                  \begin{lstlisting}[language=Java]
// 源码
interface Runnable { void run(); }

// 字节码
// flags: (0x0600) ACC_INTERFACE, ACC_ABSTRACT
// interface标记为接口类型
				      \end{lstlisting}
				            \item \textbf{代码示例}：
				                  \begin{lstlisting}[language=Java]
// 定义接口（不是继承）
interface Cloneable {
    void clone();
}

// 实现接口（不是继承）
class MyClass implements Cloneable {
    public void clone() {}
}
				      \end{lstlisting}
			            \end{itemize}

			      \item[D] \textbf{abstract — 错误（抽象修饰符）}
			            \begin{itemize}[label={}, leftmargin=*, nosep]
				            \item \textbf{语法规则}：abstract class / abstract method
				            \item \textbf{作用}：修饰符，不是实现语法
				            \item \textbf{字节码标记}：
				                  \begin{lstlisting}[language=Java]
// 源码
abstract class Shape {
    abstract void draw();
}

// 字节码
// flags: (0x0400) ACC_ABSTRACT
// ACC_ABSTRACT标记抽象类和方法
				      \end{lstlisting}
				            \item \textbf{完整示例}：
				                  \begin{lstlisting}[language=Java]
// abstract修饰抽象类
abstract class Animal {
    String name;
    Animal(String name) {
        this.name = name;
    }

    // abstract修饰抽象方法
    abstract void makeSound();
}

// 子类继承抽象类，实现抽象方法
class Cat extends Animal {
    Cat(String name) {
        super(name);
    }

    @Override
    void makeSound() {
        System.out.println(name + " 喵喵叫");
    }
}
// abstract是修饰符，extends是实现继承的关键字
				      \end{lstlisting}
			            \end{itemize}
		      \end{itemize}

		      \textbf{【应背诵知识点】}
		      \begin{itemize}[label={}, nosep]
			      \item 类继承：class 子类 extends 父类（单继承）
			      \item 接口实现：class 类 implements 接口（多实现）
			      \item 选项辨析：
			            \begin{itemize}[label={}, nosep]
				            \item extends→继承类
				            \item implements→实现接口
				            \item interface→定义接口
				            \item abstract→抽象修饰符
			            \end{itemize}
		      \end{itemize}

		      \textbf{【命题趋势与实战策略】}

		      \textbf{考场秒杀三步法}：
		      \begin{enumerate}[label={}, nosep]
			      \item 看到"继承类"→extends
			      \item 看到"实现接口"→implements
			      \item 看到"实现继承"→找extends
		      \end{enumerate}

		      \textbf{记忆口诀}：
		      \begin{itemize}[label={}, nosep]
			      \item 类继承用extends，单根继承是规则
			      \item 接口实现implements，多实现也允许
		      \end{itemize}

		      \textbf{高频陷阱}：
		      \begin{enumerate}[label={}, nosep]
			      \item 混淆extends和implements
			      \item 误认为interface是继承关键字
			      \item 误认为abstract可替代extends
		      \end{enumerate}

		      \textbf{5秒决策技巧}：看到"实现继承"→extends
	      \end{bluebox}

	      \vspace{1em}

	\item \textbf{以下关于抽象类的说法错误的是}
	      \begin{enumerate}[label=\Alph*.]
		      \item 抽象类可以有抽象方法
		      \item 抽象类不能被实例化
		      \item 抽象类的子类必须实现所有抽象方法
		      \item 抽象类中可以有非抽象方法
	      \end{enumerate}

	      \begin{bluebox}{答案与深度解析}
		      \textbf{答案：C. 抽象类的子类必须实现所有抽象方法}

		      \textbf{【核心认知框架】}
		      \begin{itemize}[label={}, leftmargin=*, nosep]
			      \item \textbf{题目本质}：考查抽象类与子类的实现规则
			      \item \textbf{关键区分点}：抽象子类可延迟实现，具体子类必须实现
			      \item \textbf{命题人意图}：测试抽象类的灵活使用规则
		      \end{itemize}

		      \textbf{【选项原理深度拆解】}

		      \begin{itemize}[leftmargin=*, nosep, label={}]
			      \item[A] \textbf{抽象类可以有抽象方法 — 正确}
			            \begin{itemize}[label={}, leftmargin=*, nosep]
				            \item \textbf{JVM规范}：抽象类可包含抽象方法（JLS §8.1.1.1）
				            \item \textbf{字节码证据}：
				                  \begin{lstlisting}[language=Java]
// 源码
abstract class Shape {
    abstract void draw();  // 抽象方法
}

// 字节码
// flags: (0x0400) ACC_ABSTRACT
// abstract void draw();  // ACC_ABSTRACT修饰方法
				            \end{lstlisting}
				            \item \textbf{代码示例}：
				                  \begin{lstlisting}[language=Java]
abstract class Animal {
    String name;

    // 抽象方法（无方法体）
    abstract void makeSound();

    // 非抽象方法
    void sleep() {
        System.out.println(name + " 睡觉");
    }
}
// 抽象类可以有0个或多个抽象方法
// 必须包含至少一个抽象方法才值得使用abstract关键字
				      \end{lstlisting}
			            \end{itemize}

			      \item[B] \textbf{抽象类不能被实例化 — 正确}
			            \begin{itemize}[label={}, leftmargin=*, nosep]
				            \item \textbf{JVM机制}：new指令不能实例化抽象类（JLS §15.9.1）
				            \item \textbf{编译器检查}：
				                  \begin{lstlisting}[language=Java]
abstract class Shape {
    abstract void draw();
}

// 编译错误！
Shape shape = new Shape();
// error: Shape is abstract; cannot be instantiated
				      \end{lstlisting}
				            \item \textbf{字节码验证}：
				                  \begin{lstlisting}[language=Java]
// 尝试实例化抽象类时编译失败
// javac编译时：
// - 检查ACC_ABSTRACT标志
// - 如果flag=ACC_ABSTRACT则拒绝new操作
// 输出编译错误信息
				            \end{lstlisting}
				            \item \textbf{正确用法}：
				                  \begin{lstlisting}[language=Java]
// 正确：通过子类实例化
abstract class Shape {
    abstract void draw();
}

class Circle extends Shape {
    @Override
    void draw() {
        System.out.println("画圆形");
    }
}

Shape shape = new Circle();  // 实例化子类（向上转型）
shape.draw();  // 输出：画圆形
				      \end{lstlisting}
				            \item \textbf{匿名内部类实例化}：
				                  \begin{lstlisting}[language=Java]
// 使用匿名内部类"实现"抽象类
Shape shape = new Shape() {
    @Override
    void draw() {
        System.out.println("画图形");
    }
};
shape.draw();  // 输出：画图形
// 本质：创建子类并实例化
				      \end{lstlisting}
			            \end{itemize}

			      \item[C] \textbf{抽象类的子类必须实现所有抽象方法 — 错误（本题答案）}
			            \begin{itemize}[label={}, leftmargin=*, nosep]
				            \item \textbf{错误原因}：子类若也是抽象类，可不实现
				            \item \textbf{正确规则}：
				                  \begin{itemize}[label={}, nosep]
					                  \item 具体子类：必须实现所有抽象方法
					                  \item 抽象子类：可不实现（可延迟到更底层）
				                  \end{itemize}
				            \item \textbf{代码证据}：
				                  \begin{lstlisting}[language=Java]
abstract class Animal {
    abstract void makeSound();
    abstract void move();
}

// 抽象子类：不实现所有抽象方法
abstract class Dog extends Animal {
    // 只实现一个抽象方法
    @Override
    void makeSound() {
        System.out.println("汪汪叫");
    }
    // move()未实现，Dog必须声明为abstract

    abstract void bark();  // 可以新增抽象方法
}

// 具体子类：必须实现所有抽象方法
class Poodle extends Dog {
    // 必须实现Dog的抽象方法bark()
    // 必须实现Animal遗留的抽象方法move()
    @Override
    void bark() {
        System.out.println("贵宾犬吠叫");
    }

    @Override
    void move() {
        System.out.println("贵宾犬跑步");
    }
}
// 抽象子类可延迟实现具体子类必须实现
				      \end{lstlisting}
				            \item \textbf{命题陷阱}：考生误认为"子类=具体子类"，忽略了抽象子类
				            \item \textbf{字节码标记}：
				                  \begin{lstlisting}[language=Java]
// 抽象子类的字节码
// flags: (0x0400) ACC_ABSTRACT
// 即使有未实现的抽象方法，只要标记为abstract即可
				            \end{lstlisting}
			            \end{itemize}

			      \item[D] \textbf{抽象类中可以有非抽象方法 — 正确}
			            \begin{itemize}[label={}, leftmargin=*, nosep]
				            \item \textbf{JVM规范}：抽象类可包含非抽象方法（JLS §8.4）
				            \item \textbf{字节码证据}：
				                  \begin{lstlisting}[language=Java]
// 源码
abstract class Shape {
    void common() {  // 非抽象方法
        System.out.println("通用方法");
    }
    abstract void draw();
}

// 字节码：
// 方法表包含两个方法
// 0: common()  // 普通方法
// 1: draw()    // ACC_ABSTRACT标记
				            \end{lstlisting}
				            \item \textbf{代码示例}：
				                  \begin{lstlisting}[language=Java]
abstract class Animal {
    String name;

    // 构造方法（非抽象）
    Animal(String name) {
        this.name = name;
    }

    // 实例方法（非抽象）
    void sleep() {
        System.out.println(name + " 睡觉");
    }

    // 静态方法（非抽象）
    static void info() {
        System.out.println("动物信息");
    }

    // 抽象方法
    abstract void makeSound();
}

// 子类继承所有非抽象方法
class Cat extends Animal {
    Cat(String name) {
        super(name);
    }

    @Override
    void makeSound() {
        System.out.println(name + " 喵喵叫");
        sleep();  // 调用继承的非抽象方法
    }
}
// 抽象类可包含构造方法、实例方法、静态方法
				      \end{lstlisting}
				            \item \textbf{模板方法模式}：
				                  \begin{lstlisting}[language=Java]
abstract class DataProcessor {
    // 模板方法（非抽象）
    public final void process() {
        loadData();       // 具体
        processData();    // 抽象（子类实现）
        saveData();       // 具体
    }

    private void loadData() { /* 加载数据 */ }
    abstract void processData();  // 子类实现
    private void saveData() { /* 保存数据 */ }
}

class CSVProcessor extends DataProcessor {
    @Override
    void processData() {
        System.out.println("处理CSV数据");
    }
}
// 抽象类定义算法骨架（非抽象），子类填充细节
				      \end{lstlisting}
			            \end{itemize}
		      \end{itemize}

		      \textbf{【应背诵知识点】}
		      \begin{itemize}[label={}, nosep]
			      \item 抽象类特性：不能实例化、可有/无抽象方法、可有非抽象方法
			      \item 子类规则：
			            \begin{itemize}[label={}, nosep]
				            \item 具体子类→必须实现所有抽象方法
				            \item 抽象子类→可不实现（延迟到下级）
			            \end{itemize}
			      \item 选项辨析：
			            \begin{itemize}[label={}, nosep]
				            \item 有抽象方法→正确
				            \item 不能实例化→正确
				            \item 子类必须实现所有→错误（抽象子类可不实现）
				            \item 有非抽象方法→正确
			            \end{itemize}
		      \end{itemize}

		      \textbf{【命题趋势与实战策略】}

		      \textbf{考场秒杀三步法}：
		      \begin{enumerate}[label={}, nosep]
			      \item 看到"必须所有"→警惕（可能是陷阱）
			      \item 区分"抽象子类"和"具体子类"
			      \item 子类是abstract→可不实现
		      \end{enumerate}

		      \textbf{记忆口诀}：
		      \begin{itemize}[label={}, nosep]
			      \item 抽象类不能new，方法可有也有无
			      \item 抽象子类可延后，具体子类要全做
		      \end{itemize}

		      \textbf{高频陷阱}：
		      \begin{enumerate}[label={}, nosep]
			      \item 误认为"所有子类必须实现所有抽象方法"
			      \item 忽略抽象子类可不实现的事实
			      \item 混淆"实例化"和"作为类型引用"
		      \end{enumerate}

		      \textbf{5秒决策技巧}：看到"子类必须实现所有抽象方法"→错误（抽象子类例外）
	      \end{bluebox}

	      \vspace{1em}

	\item \textbf{以下哪个关键字用于定义接口？}
	      \begin{enumerate}[label=\Alph*.]
		      \item interface
		      \item abstract class
		      \item class
		      \item extends
	      \end{enumerate}

	      \begin{bluebox}{答案与深度解析}
		      \textbf{答案：A. interface}

		      \textbf{【核心认知框架】}
		      \begin{itemize}[label={}, leftmargin=*, nosep]
			      \item \textbf{题目本质}：考查接口定义的关键字
			      \item \textbf{关键区分点}：interface vs abstract class vs class
		      \end{itemize}

		      \textbf{【选项原理深度拆解】}

		      \begin{itemize}[leftmargin=*, nosep, label={}]
			      \item[A] \textbf{interface — 正确（接口定义关键字）}
			            \begin{itemize}[label={}, leftmargin=*, nosep]
				            \item \textbf{语法规则}：interface 接口名（JLS §9.1）
				            \item \textbf{字节码证据}：
				                  \begin{lstlisting}[language=Java]
// 源码
interface Flyable { void fly(); }

// 字节码（javap -v）
// flags: (0x0600) ACC_INTERFACE, ACC_ABSTRACT
// Access flags: 0x0600
// 代表接口类型
				            \end{lstlisting}
				            \item \textbf{代码示例}：
				                  \begin{lstlisting}[language=Java]
// 定义接口
interface Flyable {
    void fly();  // 隐式public abstract
    int MAX_ALTITUDE = 10000;  // 隐式public static final
}

// 实现接口
class Bird implements Flyable {
    @Override
    public void fly() {  // 必须public
        System.out.println("飞翔");
    }
}
				      \end{lstlisting}
				            \item \textbf{接口特性}：
				                  \begin{itemize}[label={}, nosep]
					                  \item 方法默认public abstract
					                  \item 字段默认public static final
					                  \item 不能实例化
					                  \item 多接口继承
				                  \end{itemize}
			            \end{itemize}

			      \item[B] \textbf{abstract class — 错误（抽象类定义）}
			            \begin{itemize}[label={}, leftmargin=*, nosep]
				            \item \textbf{语法}：abstract class 类名
				            \item \textbf{区别}：抽象类可有构造方法、实例变量、非抽象方法
				            \item \textbf{代码示例}：
				                  \begin{lstlisting}[language=Java]
// 抽象类（不是接口）
abstract class Animal {
    String name;  // 实例变量

    Animal(String name) {  // 构造方法
        this.name = name;
    }

    void sleep() {  // 非抽象方法
        System.out.println("睡觉");
    }

    abstract void makeSound();  // 抽象方法
}
// abstract class是抽象类，不是接口
				      \end{lstlisting}
			            \end{itemize}

			      \item[C] \textbf{class — 错误（普通类定义）}
			            \begin{itemize}[label={}, nosep]
				            \item \textbf{语法}：class 类名
				            \item \textbf{区别}：class可实例化，不能定义接口
				            \item \textbf{字节码}：
				                  \begin{lstlisting}[language=Java]
// 源码
class Person {}

// 字节码
// flags: (0x0020) ACC_SUPER
// 无ACC_INTERFACE标志
				            \end{lstlisting}
			            \end{itemize}

			      \item[D] \textbf{extends — 错误（继承关键字）}
			            \begin{itemize}[label={}, nosep]
				            \item \textbf{用法}：子类 extends 父类
				            \item \textbf{区别}：extends用于继承，interface用于定义
			            \end{itemize}
		      \end{itemize}

		      \textbf{【应背诵知识点】}
		      \begin{itemize}[label={}, nosep]
			      \item 接口定义：interface 接口名
			      \item 选项辨析：
			            \begin{itemize}[label={}, nosep]
				            \item interface→定义接口
				            \item abstract class→定义抽象类
				            \item class→定义普通类
				            \item extends→继承操作符
			            \end{itemize}
		      \end{itemize}

		      \textbf{【命题趋势与实战策略】}

		      \textbf{5秒决策技巧}：看到"定义接口"→interface
	      \end{bluebox}

	\item \textbf{以下关于接口的说法错误的是}
	      \begin{enumerate}[label=\Alph*.]
		      \item 接口中的方法默认是public abstract修饰的
		      \item 接口中的变量默认是public static final修饰的
		      \item 一个类可以实现多个接口
		      \item 接口可以继承类
	      \end{enumerate}

	      \begin{bluebox}{答案与深度解析}
		      \textbf{答案：D. 接口可以继承类}

		      \textbf{【核心认知框架】}
		      \begin{itemize}[label={}, leftmargin=*, nosep]
			      \item \textbf{题目本质}：考查接口的继承规则和成员默认修饰符
			      \item \textbf{关键区分点}：接口只能继承接口，不能继承类
			      \item \textbf{命题人意图}：测试对接口机制的深度理解
		      \end{itemize}

		      \textbf{【选项原理深度拆解】}

		      \begin{itemize}[leftmargin=*, nosep, label={}]
			      \item[A] \textbf{接口中的方法默认是public abstract修饰的 — 正确}
			            \begin{itemize}[label={}, leftmargin=*, nosep]
				            \item \textbf{JLS规范}：JLS §9.4 接口方法隐式public abstract
				            \item \textbf{字节码证据}：
				                  \begin{lstlisting}[language=Java]
// 源码
interface Flyable { void fly(); }

// 字节码（javap -v）
// public abstract void fly()
// flags: ACC_PUBLIC, ACC_ABSTRACT
						      \end{lstlisting}
				            \item \textbf{代码示例}：
				                  \begin{lstlisting}[language=Java]
interface MyInterface {
    // 三种声明等价
    void method1();
    public void method2();
    abstract void method3();

    // default和static方法例外（Java 8+）
    default void defaultMethod() {
        System.out.println("默认实现");
    }
    static void staticMethod() {
        System.out.println("静态方法");
    }
}
// 普通接口方法默认public abstract
					      \end{lstlisting}
			            \end{itemize}

			      \item[B] \textbf{接口中的变量默认是public static final修饰的 — 正确}
			            \begin{itemize}[label={}, leftmargin=*, nosep]
				            \item \textbf{JLS规范}：JLS §9.3 接口字段隐式public static final
				            \item \textbf{字节码证据}：
				                  \begin{lstlisting}[language=Java]
// 源码
interface Constants { int MAX = 100; }

// 字节码
// field public static final int MAX
// flags: ACC_PUBLIC, ACC_STATIC, ACC_FINAL
// ConstantValue: int 100
						      \end{lstlisting}
				            \item \textbf{代码示例}：
				                  \begin{lstlisting}[language=Java]
interface Constants {
    // 三种声明等价
    int MAX = 100;
    public static final int MIN = 0;
    final int COUNT = 10;

    // 错误：不能是其他修饰符
    // private int x = 1;  // 编译错误
    // protected int y = 2; // 编译错误
    // int VALUE;         // 必须初始化
}

// 使用
int x = Constants.MAX;  // 通过接口名访问
// Constants.MAX = 200;  // 编译错误，final不可修改
					      \end{lstlisting}
			            \end{itemize}

			      \item[C] \textbf{一个类可以实现多个接口 — 正确}
			            \begin{itemize}[label={}, leftmargin=*, nosep]
				            \item \textbf{JLS规范}：JLS §8.1.5 多接口实现
				            \item \textbf{字节码层面}：class文件中的interfaces字段数组存储实现的接口
				            \item \textbf{完整代码示例}：
				                  \begin{lstlisting}[language=Java]
// 定义多个接口
interface Flyable { void fly(); }
interface Runnable { void run(); }
interface Swimmable { void swim(); }

// 实现多个接口
class Duck implements Flyable, Runnable, Swimmable {
    @Override
    public void fly() {
        System.out.println("鸭子飞翔");
    }

    @Override
    public void run() {
        System.out.println("鸭子跑");
    }

    @Override
    public void swim() {
        System.out.println("鸭子游泳");
    }
}

// 使用
Duck duck = new Duck();
duck.fly();   // 实现Flyable
duck.run();   // 实现Runnable
duck.swim();  // 实现Swimmable
// 一个类可实现多个接口，弥补单继承局限
					      \end{lstlisting}
				            \item \textbf{默认方法冲突处理}：
				                  \begin{lstlisting}[language=Java]
interface A {
    default void method() {
        System.out.println("A.method");
    }
}

interface B {
    default void method() {
        System.out.println("B.method");
    }
}

// 编译错误：多个默认方法冲突
class C implements A, B {}  // error: 接口A和B都继承了method()

// 解决：显式重写并指定调用
class D implements A, B {
    @Override
    public void method() {
        A.super.method();  // 或 B.super.method()
    }
}
// 多接口实现的默认方法冲突必须显式解决
					      \end{lstlisting}
			            \end{itemize}

			      \item[D] \textbf{接口可以继承类 — 错误（本题答案）}
			            \begin{itemize}[label={}, leftmargin=*, nosep]
				            \item \textbf{错误原因}：接口只能继承接口，不能继承类
				            \item \textbf{JLS规范}：JLS §9.1 接口extends子句只能列出接口类型
				            \item \textbf{底层原因}：
				                  \begin{itemize}[label={}, nosep]
					                  \item 接口是规范，不包含实现代码
					                  \item 类包含实现，违反接口设计原则
					                  \item JVM不支持接口继承类的字节码结构
					                  \item 破坏Java的单继承模型
				                  \end{itemize}
				            \item \textbf{字节码证据}：
				                  \begin{lstlisting}[language=Java]
// 编译失败的代码
class MyClass {}
interface MyInterface extends MyClass {}  // 编译错误

// javac错误信息：
// expected: 接口类型
// found: 类MyClass
// 接口只能继承其他接口
						      \end{lstlisting}
				            \item \textbf{代码反证}：
				                  \begin{lstlisting}[language=Java]
class Animal {
    void eat() {}  // 实现方法
}

// 编译错误！接口不能继承类
interface Pet extends Animal {
    // ^ 提示：期望接口类型
}

// 正确做法：接口继承接口
interface AnimalBehavior {
    void eat();
    void sleep();
}

interface DogBehavior extends AnimalBehavior {
    void bark();
}
// 接口只能extends接口，不能extends类
					      \end{lstlisting}
				            \item \textbf{接口的多继承}：
				                  \begin{lstlisting}[language=Java]
interface A { void a(); }
interface B { void b(); }
interface C { void c(); }

// 接口可以继承多个接口（接口的多继承）
interface All extends A, B, C {
    // All拥有a()、b()、c()三个抽象方法
}

class Impl implements All {
    public void a() {}
    public void b() {}
    public void c() {}
}
// 接口多重继承是允许的，但只能继承接口
					      \end{lstlisting}
			            \end{itemize}
		      \end{itemize}

		      \textbf{【接口vs抽象类继承对比】}
		      \begin{itemize}[label={}, nosep]
			      \item 类继承：class Son extends Father（单继承一个类）
			      \item 类实现：class Son implements I1, I2（多实现多个接口）
			      \item 接口继承：interface I extends J, K（多继承接口，不能继承类）
			      \item 接口继承只能指向接口，不能指向类
		      \end{itemize}

		      \textbf{【应背诵知识点】}
		      \begin{itemize}[label={}, nosep]
			      \item 接口方法隐式public abstract
			      \item 接口字段隐式public static final
			      \item 类可多实现接口
			      \item 接口只能继承接口
			      \item 选项辨析：
			            \begin{itemize}[label={}, nosep]
				            \item 方法默认pub抽象→正确
				            \item 字段默认pub静final→正确
				            \item 类可多实现→正确
				            \item 接口能继承类→错误
			            \end{itemize}
		      \end{itemize}

		      \textbf{【命题趋势与实战策略】}

		      \textbf{考场秒杀三步法}：
		      \begin{enumerate}[label={}, nosep]
			      \item 看到方法→默认public abstract
			      \item 看到变量→默认public static final
			      \item 看到接口继承类→直接否定
		      \end{enumerate}

		      \textbf{记忆口诀}：
		      \begin{itemize}[label={}, nosep]
			      \item 接口方法pub抽象，字段pub静final
			      \item 类可多实现接口
			      \item 接口只能继接口，不能继承类
		      \end{itemize}

		      \textbf{高频陷阱}：
		      \begin{enumerate}[label={}, nosep]
			      \item 误认为接口可以继承类
			      \item 混淆接口多继承（接口）和类多实现（接口实现）
		      \end{enumerate}

		      \textbf{5秒决策技巧}：看到"接口继承类"→错
	      \end{bluebox}

\end{enumerate}

\end{document}
